\section{工作流与事故演练}

\subsection{部署前演练}

每次 deployment 之前，应模拟一轮典型任务：从 dataset、signature、optimizer 到 evidence 结果。演练的目标是确保 DevOps 团队与 optimizer 得出的 prompt 保持一致。

演练流程最好在 staging 环境中完成：把 evidence snapshot 上传至 dashboard，让 QA 团队先验算每个 module 的 metric。只有当 evidence run 与 production run 相符（误差 < 5%），才允许上线。

\subsection{事故响应流程}

事故响应流程包括三步：1) trace → 2) rollback → 3) annotate。Trace 阶段要快速定位哪个 module 产生 drift；rollback 阶段要锁定 last-known-good prompt；annotate 阶段要把事件写入 governance log。

回放 trace 时要附上 timeline（Action Timeline / Evidence Matrix），并对比 metric drift 时间戳。这能帮助团队理解 drift 是 optimizer 输入还是 tool output 侧的问题。

\begin{importantbox}{事故演练 checklist}
\begin{itemize}
    \item 确认 alarm 触发条件（latency/metric/feedback drift）
    \item 暂停 optimizer，并保留 evidence snapshot
    \item 将 recovery steps 记录在 Action Timeline，包含 owner + timestamp
\end{itemize}
\end{importantbox}

\subsection{本章小结}

工作流与事故演练把 DSPy pipeline 从 research 级别推向 production，通过 checklist 和 Action Timeline 让团队能在事件发生时快速恢复。

